% answers_short.tex —— 2010 年·普通高等学校招生考试·湖南卷·文科 · 简短答案速查
% 由 tools/make_short_answers.py 自 answers.tex 的 \gksolhead/\gkanswer 提取，
% 并与 2010/2010_数学真题_湖南文_答案要点.md 逐题比对后生成，勿手改（改源 MD 后重跑）。

\begin{gknotice}
\textbf{使用说明}：本节只给最终答案便于快速核对，完整推导见解析版。答案为编写者独立推导，\textbf{非官方标准答案}。
\end{gknotice}

\gkpart{选择题}
\begin{center}\begin{tabular}{ccccccccc}
\toprule
题号 & 1 & 2 & 3 & 4 & 5 & 6 & 7 & 8 \\
\midrule
答案 & A & C & A & D & B & C & A & D \\
\bottomrule
\end{tabular}\end{center}

\gkpart{填空题}
\begin{itemize}[leftmargin=2.6em]
  \item[\textbf{9.}] $3$。
  \item[\textbf{10.}] $161.8$ g 或 $138.2$ g。
  \item[\textbf{11.}] $\dfrac13$。
  \item[\textbf{12.}] $x\geqslant0$（填 $x>0$ 亦可）。
  \item[\textbf{13.}] $4$。
  \item[\textbf{14.}] $-1$；$x^2+(y-1)^2=1$。
  \item[\textbf{15.}] $5$；$\{a_1,a_2,a_5,a_7,a_8\}$。
\end{itemize}

\gkpart{解答题}
\begin{itemize}[leftmargin=2.6em]
  \item[\textbf{16.}] （1）$\pi$；（2）最大值为 $\sqrt2-1$，此时 $x$ 的集合为 $\left\{x\ \middle|\ x=k\pi+\dfrac\pi8,\ k\in\mathbf Z\right\}$。
  \item[\textbf{17.}] （1）$x=1$，$y=3$；（2）$\dfrac3{10}$。
  \item[\textbf{18.}] （1）$\sqrt2$；（2）证明见解析。
  \item[\textbf{19.}] （1）$\dfrac{x^2}{25}+\dfrac{y^2}{9}=1$；（2）$5$ 年。
  \item[\textbf{20.}] （1）表 $4$ 各行为 $1,3,5,7$；$4,8,12$；$12,20$；$32$，各行平均数为 $4,8,16,32$，推广：表 $n$ 各行的平均数构成首项为 $n$、公比为 $2$ 的等比数列；（2）$4-\dfrac{1}{(n+1)2^{n-2}}$。
  \item[\textbf{21.}] （1）$a<-1$ 时在 $(0,1)$、$(-a,+\infty)$ 上单调递增，在 $(1,-a)$ 上单调递减；$-1<a<0$ 时在 $(0,-a)$、$(1,+\infty)$ 上单调递增，在 $(-a,1)$ 上单调递减；（2）存在，$a\in[-3,-2]$。
\end{itemize}
